\chapter*{Conclusion}
\chaptermark{Conclusion}
\markboth{CONCLUSION}{CONCLUSION}
\addcontentsline{toc}{chapter}{Conclusion}
\label{ch:Conclusion}


This thesis presented a search for Heavy Stable Charged Particles in proton--proton collisions recorded by the CMS experiment during the Run--3 data-taking period of the LHC, at a centre-of-mass energy of $\sqrt{s} = 13.6$~TeV. The work presents two complementary contributions: an instrumental study dedicated to reconstructing the ionisation energy loss (\dedx) in the CMS silicon strip tracker, and the design and validation of the analysis strategy for the search itself, interpreted in terms of gluino pair production leading to R-hadrons.

\vspace{0.5 cm}

The first contribution addresses the saturation of the strip readout electronics, which biases the \dedx measurement precisely in the regime relevant to this search: highly ionising deposits left by slow, massive particles. A new correction based on the shape of the strip clusters was developed, classifying clusters according to the position of the saturated strip within the deposit and treating configurations with one or two consecutive saturated strips. The method was validated on both simulation and 2018 collision data, and was approved by the CMS Tracking Physics Object Group. Its performance was then demonstrated on the full 2024 and 2025 datasets by measuring the average ionisation $I_h$ as a function of track momentum. The benefit for the search is direct: for a simulated gluino of mass $2000$~GeV/$c^2$, the reconstructed mass is $1688 \pm 533$~GeV/$c^2$ without any saturation correction and $1900 \pm 503$~GeV/$c^2$ with the correction used previously, while the shape correction restores it to $2053 \pm 525$~GeV/$c^2$. The correction thus removes most of the systematic bias affecting the mass reconstruction, on which both the background estimation and the signal extraction of the analysis rely.

The second contribution is the search itself, performed on the 2024 dataset corresponding to an integrated luminosity of 110~fb$^{-1}$. Two limitations of the Run--2 analysis were addressed. First, the event selection is based on missing transverse energy triggers rather than muon triggers, recovering R-hadrons that flip charge or become neutral in the detector and would otherwise escape muon-based selections; this strategy roughly triples the event selection efficiency for gluino R-hadrons, from about 11\% to about 30\%. Second, the acceptance was extended to the full pseudo-rapidity coverage of the tracker, with the barrel ($|\eta| < 1$) and endcap ($1 \leq |\eta| < 2.4$) regions treated as separate channels. Candidate masses are reconstructed by combining the momentum and the $I_h$ estimator, and the background is predicted directly from data with an ABCD method exploiting the independence of the pixel discriminator $F_{\text{pixel}}$ and the transverse momentum for Standard Model tracks. The mass reconstruction and the background prediction were validated both in data, in dedicated validation regions, and in simulation, and the associated systematic uncertainties were derived. As the signal region remains blinded, the results are reported as expected limits. In the most sensitive configuration, combining the barrel and endcap channels with a shape-based extraction, gluino masses up to 2.42~TeV are expected to be excluded at 95\% confidence level, to be compared with the 2.08~TeV observed exclusion of the Run--2 search at a similar integrated luminosity. The expected discovery significance exceeds $5\sigma$ over most of the accessible mass range and remains close to $3\sigma$ near 2600~GeV, so that the analysis retains a genuine discovery potential in the region beyond the current limits.

Several steps remain before unblinding. The residual population of high-mass events observed in the validation region, dominated by $t\bar{t}$ production, is being characterised, and a dedicated selection criterion may be added to the pre-selection, benefiting mainly the endcap channel where the signal-to-background ratio is lower. The statistical treatment will also be refined: the limits presented here rely on the asymptotic approximation, which is known to be optimistic in the low-count regime of the barrel channel, and the final results will be recomputed with the full CL$_\text{s}$ construction based on pseudo-experiments. Once unblinded, the analysis will directly probe the mass window in which the ATLAS collaboration reported a $3.6\sigma$ excess in its full Run--2 search, and is therefore in a position to confirm or refute this observation with an independent detector and methodology.

Beyond the present iteration, the sensitivity of the search is limited by the expected signal yield rather than by systematic uncertainties, and the most significant gain will come from the full Run--3 dataset, roughly three times larger than the one analysed here. A simple scaling argument places the corresponding mass reach at 2600--2700~GeV. The analysis chain is also directly applicable to R-hadrons formed from pair-produced top squarks, requiring only the processing of the corresponding signal samples, while interpretations for finite lifetimes and charge-suppression scenarios would further enlarge the covered phase space. Finally, combining the missing transverse energy triggers with the muon paths in a logical OR would recover lepton-like candidates such as tau sleptons within a single analysis framework.

On a longer timescale, the Phase--2 upgrade of CMS will profoundly modify the experimental inputs of this search. The binary readout of the outer tracker will restrict the $dE/dx$ measurement to the upgraded inner tracker, reducing the average number of ionisation samples per track from about fifteen to at most six, while the MIP Timing Detector will provide time-of-flight measurements with a resolution of 30--60~ps, making timing the primary signature for slow particles at the HL--LHC. The completion of this search on the full Run--3 dataset therefore represents both the culmination of the ionisation-based strategy pursued by CMS since Run--1 and the most sensitive test of the R-hadron hypothesis achievable before the advent of timing-based methods.